\subsection*{先问：我们究竟想记录什么？}
假设你在地图上先向东走两格，再向南走一格。别人从另一个地点照做，到达的地点不同，但完成了同一个位移。向量首先记录这种可以搬动而不改变的动作。画箭头时，起点可以换，方向与长度却不能随意换。“从原点出发的箭头”只是方便的画法，并不是向量的全部含义。

把这段路写成 $(2,-1)$，两个数其实是两条指令：沿第一个基方向走两份，沿第二个基方向走负一份。坐标不是贴在箭头上的永久标签；换一套方向与单位，同一位移会得到另一组数。因此阅读坐标之前，要先问坐标轴怎样选、每格代表什么。若两个分量分别记录温度和质量，也不能直接赋予欧氏长度一个物理意义。

\subsection*{让运算来自动作}
再接着走 $(-1,3)$，横向净移动一格、纵向净移动两格，所以相加可以逐坐标进行。把整条路线走两遍，每一步都扩大两倍；“先合并再扩大”与“分别扩大再合并”得到同一结果。分配律由此成为两种操作顺序的一致性要求。

先画首尾相接的箭头，再把它们搬到共同起点，补成平行四边形。两幅图为什么指向同一终点？回答这个问题，就把几何相加与坐标相加接了起来。后面遇到有几百个分量的数据向量，虽然画不出箭头，逐项加法与整体缩放仍保留这套结构。下一节将追问：只给少数几个方向，能拼出多少向量？
